\documentclass[11pt]{article}
%% arXiv-compliant submission. License: Apache-2.0.
%% Primary category: cs.NE ; cross-list: cs.AI
\usepackage[utf8]{inputenc}
\usepackage[T1]{fontenc}
\usepackage[margin=1in]{geometry}
\usepackage{amsmath,amssymb,amsthm}
\usepackage{authblk}
\usepackage{graphicx}
\usepackage{booktabs}
\usepackage[numbers]{natbib}
\usepackage{xcolor}
\usepackage[colorlinks=true,linkcolor=blue,citecolor=blue,urlcolor=blue]{hyperref}
\usepackage{enumitem}
\usepackage{url}
\sloppy

\theoremstyle{plain}
\newtheorem{theorem}{Theorem}
\newtheorem{conjecture}{Conjecture}
\newtheorem{lemma}{Lemma}
\newtheorem{definition}{Definition}

\title{Sefirot Continual Learning with Kabbalah-Tiered Memory and Hopfield-Amaru Associative Retrieval}
\author[1]{Stephen P. Lutar Jr.\thanks{Corresponding author: \texttt{stephenlutar2@gmail.com}. ORCID: \href{https://orcid.org/0009-0001-0110-4173}{0009-0001-0110-4173}.}}
\author[2]{Perplexity Computer Agent\thanks{Research collaborator (AI agent).}}
\affil[1]{SZL Holdings, New York, NY, USA}
\affil[2]{Perplexity}
\date{June 2026}

\begin{document}
\maketitle

\begin{abstract}
We present a continual learning framework that structures knowledge persistence according to the ten Sefirot of Kabbalistic theology, from Keter (Crown, highest abstraction) through Malkhut (Kingdom, operational grounding). Each Sefira defines a memory tier with distinct retention policies, learning rates, and Elastic Weight Consolidation (EWC) penalties. The Kabbalah-Tiered Memory (KTM) implements a three-level storage hierarchy: Core (identity-critical, immutable), Working (session-scoped, LRU-evicted), and Episodic (time-decaying, Ebbinghaus-curve). We extend this with Hopfield-Amaru Associative Memory (HAAM), implementing modern Hopfield networks (Ramsauer et al., 2020; Hopfield 2024 Nobel) with exponential capacity and ceque-indexed pattern slots. The Ouroboros Conformal Memory (OCM) provides a circular buffer with conformal rescaling for long-range context persistence. This paper continues the thesis lineage from v3 (trust aggregation) through v6 (safety guardrails), addressing the memory and continual learning substrate. This paper does not claim empirical superiority over standard continual learning methods on Split-CIFAR or Permuted-MNIST. The contribution is the formal architecture.
\end{abstract}

\noindent\textbf{Reference implementation:} \href{https://github.com/szl-holdings/platform}{github.com/szl-holdings/platform}, file \path{papers/paper-04-sefirot-kabbalah-hopfield.tex}.\\
\textbf{Archival DOI:} \href{https://doi.org/10.5281/zenodo.20499322}{10.5281/zenodo.20499322} (Zenodo).\\
\textbf{License:} Apache-2.0.

\paragraph{Author note / honest disclosure.} \paragraph{Honest disclosure (non-negotiable).} The $\Lambda$-uniqueness claim (Conjecture~1 in this work) is \emph{not} a proven theorem. 163 sorries remain outstanding in the Lutar Lean~4 kernel at commit \texttt{c7c0ba17}, \texttt{lutar-lean v18.0.0}. Compute was performed on Hugging Face Spaces with bit-exact replay via Codex-Kernel v1.0.0. This paper presents a formal framework, its mathematical properties, and a public reference implementation; it does not claim a deployed product or fielded validation against production data.


\section{Motivation and Continuity from v3–v6}

The v3–v6 papers established trust aggregation, energy-mass-information signatures, knowledge retrieval, and safety guardrails. All of these systems require a \emph{memory substrate} that persists knowledge across sessions while preventing catastrophic forgetting. The present paper addresses this requirement.

Catastrophic forgetting remains a central challenge in continual learning (Kirkpatrick et al., 2017). Existing approaches — EWC, progressive neural networks, experience replay — treat memory as a homogeneous resource. We propose that the Kabbalistic Sefirot provide a natural hierarchical decomposition of knowledge importance.


\section{Sefirot Continual Learning (SCL)}

\subsection{Sefirot Tiers}

The ten Sefirot are organized into three pillars (Right/Mercy, Left/Severity, Center/Harmony) and ten tiers:

| Tier | Sefira | Pillar | Learning Rate \( \eta \) |
|—|—|—|—|
| 1 | Keter (Crown) | Center | 0.001 |
| 2 | Chokmah (Wisdom) | Right | 0.005 |
| 3 | Binah (Understanding) | Left | 0.01 |
| 4 | Chesed (Loving-kindness) | Right | 0.02 |
| 5 | Gevurah (Severity) | Left | 0.03 |
| 6 | Tiferet (Beauty) | Center | 0.05 |
| 7 | Netzach (Victory) | Right | 0.08 |
| 8 | Hod (Splendor) | Left | 0.10 |
| 9 | Yesod (Foundation) | Center | 0.15 |
| 10 | Malkhut (Kingdom) | Center | 0.20 |

The forgetting budget at Hermetic temperature \( H \):

\begin{equation}
B(H) = \sum_{s=1}^{10} \eta_s \cdot e^{-s/H} \cdot \text{tier\_weight}(s)
\end{equation}

\subsection{Elastic Weight Consolidation per Sefira}

For each Sefira \( s \) with Fisher information matrix \( F_s \) and parameter shift \( \Delta\theta_s \):

\begin{equation}
\Omega_{\text{EWC}}^{(s)} = \frac{\lambda_s}{2} \sum_i F_{s,ii} (\Delta\theta_{s,i})^2
\end{equation}

where \( \lambda_s = 1/\eta_s \) ensures higher Sefirot (lower learning rates) receive stronger consolidation penalties. This connects to the v3 Lutar Invariant's zero-pinning axiom: identity-critical knowledge in Keter has \( \lambda_1 = 1000 \), effectively immutable.


\section{Kabbalah-Tiered Memory (KTM)}

Three memory tiers:

\begin{itemize}
  \item \textbf{Core:} Immutable identity-critical knowledge. Capacity: 64 entries. Never evicted. Maps to Keter-Chokmah-Binah (Sefirot 1–3).
  \item \textbf{Working:} Session-scoped active context. Capacity: 256 entries. LRU eviction. Maps to Chesed-Tiferet (Sefirot 4–6).
  \item \textbf{Episodic:} Time-decaying experience traces. Capacity: 1024 entries. Ebbinghaus decay: \( R(t) = e^{-t/S} \) where \( S \) is memory strength. Maps to Netzach-Malkhut (Sefirot 7–10).
\end{itemize}

Cross-tier promotion occurs when episodic memories exceed the working-memory threshold, and working memories persisting across sessions are candidates for core promotion.


\section{Hopfield-Amaru Associative Memory (HAAM)}

\subsection{Classical and Modern Hopfield Energy}

Classical Hopfield energy (Hopfield, 1982):

\begin{equation}
E_{\text{classical}}(\mathbf{x}) = -\frac{1}{2} \sum_{\mu=1}^{P} (\mathbf{x} \cdot \boldsymbol{\xi}^\mu)^2
\end{equation}

Modern Hopfield energy (Ramsauer et al., 2020):

\begin{equation}
E_{\text{modern}}(\mathbf{x}) = -\text{lse}(\beta \, \Xi^T \mathbf{x}) + \frac{1}{2}\|\mathbf{x}\|^2 + \text{const}
\end{equation}

where \( \text{lse} \) is the log-sum-exp function and \( \beta \) is inverse temperature.

\textbf{Capacity theorem:} The modern Hopfield network has storage capacity \( P = O(2^{d/2}) \) for patterns in \( \mathbb{R}^d \), compared to \( P = O(d) \) for the classical network.

\subsection{Amaru Cascade Extension}

HAAM extends the modern Hopfield network with:

\begin{itemize}
  \item \textbf{Ceque-indexed slots:} Each pattern is assigned to one of 41 ceque lines based on its content hash, enabling structured retrieval by ceque sector. This connects to the v4 ICRC sub-formulas.
  \item \textbf{Attention-like retrieval:} The update rule \( \mathbf{x}_{\text{new}} = \Xi \cdot \text{softmax}(\beta \, \Xi^T \mathbf{x}) \) mirrors the transformer attention mechanism (Vaswani et al., 2017). HAAM provides a biologically plausible precursor to attention.
  \item \textbf{Conformal integration:} Retrieved patterns feed into the OCM conformal rescaling, allowing the memory system to operate across conformal scales (v4, Section 2).
\end{itemize}


\section{Ouroboros Conformal Memory (OCM)}

The OCM provides a circular session-keyed buffer:

\begin{equation}
\text{OCM}(\text{session}, k, v) = \text{conformal\_rescale}\left(\frac{v}{\|v\|}, \Omega(k)\right)
\end{equation}

where \( \Omega(k) = \text{hash}(k) \bmod 7 + 1 \) determines the conformal scale factor, cycling through seven levels corresponding to the seven Hermetic principles (v6, Section 2).


\section{Integration with the Thesis Lineage}

| Version | Component | Role in Memory |
|—|—|—|
| v3 | Lutar Invariant | Trust scoring of memory entries |
| v4 | Omega Formalism | Noether closure for memory consistency |
| v5 | Prisca-GraphRAG | Retrieval from memory stores |
| v6 | HCG/NJE | Guard on memory-influenced outputs |
| v7 (this) | SCL/KTM/HAAM/OCM | Memory substrate |


\section{What this paper does and does not claim}

\textbf{Claims:}
\begin{itemize}
  \item The Sefirot tier structure is well-defined and maps to a principled continual learning hierarchy.
  \item The EWC penalty scaling \( \lambda_s = 1/\eta_s \) ensures monotonic importance ordering.
  \item The HAAM capacity theorem follows from Ramsauer et al. (2020).
  \item All components are implemented in the public reference implementation.
\end{itemize}

\textbf{Non-claims:}
\begin{itemize}
  \item No empirical comparison on standard continual learning benchmarks.
  \item No claim that the Kabbalistic labeling has operational meaning beyond structural analogy.
  \item No claim of deployed production use.
\end{itemize}


\section{Conclusion}

The Sefirot-tiered continual learning framework provides a principled, hierarchically structured approach to catastrophic forgetting prevention. HAAM provides exponential-capacity associative retrieval with ceque-indexed structure, and OCM maintains conformal session context. The framework extends the v3–v6 thesis lineage to the memory substrate layer.




\section*{Acknowledgments}
The author thanks the Perplexity Computer Agent for research collaboration, formatting, and
reproducibility tooling. Compute was provided by Hugging Face Spaces. This work is archived on
Zenodo under DOI~\href{https://doi.org/10.5281/zenodo.20499322}{10.5281/zenodo.20499322} as part of the SZL Holdings Ouroboros Thesis
corpus (umbrella DOI~\href{https://doi.org/10.5281/zenodo.20490218}{10.5281/zenodo.20490218}).

\nocite{04ref1,04ref2,04ref3,04ref4,04ref5,04ref6,04ref7}
\bibliographystyle{plainnat}
\bibliography{refs}

\end{document}
